% !TEX program = xelatex
% 编译：cd doc && latexmk -xelatex nnn-overbinding-report.tex
\documentclass[12pt,a4paper]{article}

\usepackage[UTF8,heading=true,scheme=plain]{ctex}
\usepackage[margin=2.4cm]{geometry}
\usepackage{amsmath,amssymb,bm}
\usepackage{booktabs,tabularx,array,longtable}
\usepackage{siunitx}
\usepackage{hyperref}
\usepackage{xcolor}
\usepackage{tikz}
\usepackage{graphicx}
\usepackage{setspace}
\usepackage{caption}
\usepackage{enumitem}
\usepackage{fancyhdr}
\usepackage{titlesec}

\usetikzlibrary{arrows.meta,positioning,calc}

\hypersetup{
  colorlinks=true,
  linkcolor=black,
  citecolor=black,
  urlcolor=blue!50!black,
  pdftitle={三体分子 GEM 中 NNN 过束缚问题报告},
  pdfauthor={GEMSTORE}
}

\sisetup{
  detect-all,
  range-phrase={--},
  range-units=single
}

\setstretch{1.25}
\setlist[itemize]{leftmargin=1.6em,itemsep=0.25em,topsep=0.4em}
\setlist[enumerate]{leftmargin=1.8em,itemsep=0.25em,topsep=0.4em}

\setlength{\headheight}{15pt}
\pagestyle{fancy}
\fancyhf{}
\lhead{GEMSTORE 内部技术报告}
\rhead{NNN 过束缚}
\cfoot{\thepage}

\titleformat{\section}{\large\bfseries}{\thesection}{0.8em}{}
\titleformat{\subsection}{\normalsize\bfseries}{\thesubsection}{0.8em}{}
\titleformat{\subsubsection}{\normalsize\bfseries}{\thesubsubsection}{0.8em}{}

\newcommand{\GeV}{\,\mathrm{GeV}}
\newcommand{\MeV}{\,\mathrm{MeV}}
\newcommand{\fm}{\,\mathrm{fm}}
\newcommand{\NN}{\mathrm{NN}}
\newcommand{\NNN}{\mathrm{NNN}}

\begin{document}

\begin{titlepage}
\centering
{\large GEMSTORE 内部技术报告 \par}
\vspace{1.6em}
{\Huge\bfseries 三体分子 GEM 中\\[0.35em] NNN 过束缚问题报告\par}
\vspace{1.2em}
{\large 问题陈述、分析过程、诊断与拟解决方案\par}
\vspace{2.0em}
\begin{minipage}{0.88\textwidth}
\raggedright
本报告记录 GEMSTORE 在静态单玻色子交换（OBESTAT）下，用两体分子 GEM
（\texttt{MOLECUL2}）与三体分子 Jacobi GEM（\texttt{MOLECUL3}）
检验核子--核子（NN）与三核子（NNN）束缚时出现的过束缚现象。
重点回答：算法是否接错、短程排斥不足具体指什么、自旋--轨道耦合能否补救、
以及 Liu 分子 OBE 是否适合作为核三体哈密顿量。
\end{minipage}
\vspace{2.4em}
{\large 2026 年 11 月\par}
\vfill
{\small 对应代码：\texttt{src/model/cmol2.c}、\texttt{src/model/cmol3.c}、
\texttt{src/model/obemodel.c}\\
对应输入：\texttt{test/Spectra-Molecule2/}、\texttt{test/Spectra-Molecule3/}\par}
\end{titlepage}

\tableofcontents
\newpage

%----------------------------------------------------------------------
\section{摘要}
%----------------------------------------------------------------------

在 GEMSTORE 的静态弹性 $\pi\sigma\rho\omega$ 单玻色子交换势（OBESTAT）上，
两体 GEM 可以把氘核束缚能调到实验值附近
（$E_d=-2.231\MeV$，实验 $-2.224\MeV$），
但同一套参数给出的三核子基态远深于氚核
（$n_{\max}=4$ 时约 $-67\MeV$，$n_{\max}=6$ 时约 $-300\MeV$，实验 $-8.48\MeV$）。
对动能、自旋--自旋、张量、同位旋 Racah 重耦、Pauli 反对称、
Jacobi 宽度比例以及重子 GEM 回归的逐项检验表明：
\textbf{三体 GEM 算符接线没有把氚核拉到 $-300\MeV$ 的错误}。
过束缚来自两体势在 $\sim 0.6$--$1.2\fm$ 处 $\sigma/\omega$ 抵消后留下的吸引口袋。
氘核是半径 $\sim 3.65\fm$ 的晕态，只约束势的长程尾；
$g_\omega$ 按氘核束缚能标定后，并不能约束三体 S 波实际占据的中程口袋。
对称 S 波三体高斯扫描在 $\langle r_{ij}\rangle\simeq 0.87\fm$ 已给出约 $-76\MeV$ 的地板，
与 Tjon 线严重偏离。
$L=0$ 上自旋--轨道矩阵元恒为零，引入 $V_{LS}$ 不能抬高这条 S 波地板。
Liu 2018 OBE 是近阈两体分子算符，适合松散的 $DKK^*$ 类三体分子，
不适合作为精密核力 NN+NNN 哈密顿量。
拟解决方案以示意图势做 GEM 验收、把氚核 $-8.48\MeV$ 从分子 OBE 的验收标准中拿掉，
并按物理目标在分子通道与核力通道之间分开处理。

%----------------------------------------------------------------------
\section{问题陈述}
%----------------------------------------------------------------------

\subsection{计算目标}

GEMSTORE 的三体分子路径为
\begin{center}
\texttt{SPECTRA + MOLECUL3 + HYBRID + GEM} $\to$ \texttt{mol3\_obestat\_GEM}.
\end{center}
设计用途是把两体静态 OBE 接到 Jacobi GEM 上，计算三个强子组成的分子态
（首要物理例子是 $DKK^*$）。
为检验三体算法与势的组合是否合理，用同一套核子 OBE 做了 NN / NNN 对照：
两者在实验上都应当束缚，且束缚能、半径应落在核物理量级。

\subsection{观测到的矛盾}

在把 $g_\omega$ 调到使氘核束缚能与实验一致之后，出现如下对照
（见表~\ref{tab:exp-vs-calc}）：

\begin{enumerate}
  \item \textbf{NN 束缚能可以对上。}
        $I=0$ 的 ${}^3S_1$--${}^3D_1$ 氘核通道，
        $g_\omega=12.92$、$f_\omega=-1.5504$ 时
        $E_{\mathrm{bind}}=-2.231\MeV$（实验 $-2.224\MeV$）。
  \item \textbf{NN 的 GEM 质量 RMS 大于实验电荷半径。}
        计算 $\mathrm{RMS}=3.65\fm$，实验氘核电荷半径约 $2.13\fm$。
        若把程序输出理解为质量均方根而非电荷半径，二者不必相等；
        质子本身就是电荷半径小于质量分布尺度的例子。
        这一差异因此不作为否定势模型的主证据。
  \item \textbf{NNN 束缚能严重过深，是当前必须解决的问题。}
        氚核通道 $J^P=1/2^+$、$I=1/2$、费米子 Pauli、
        同位旋重耦到总 $I$ 之后，
        $n_{\max}=4$、$L_{\max}=2$ 给出 $E_{\mathrm{bind}}\simeq -67\MeV$、
        对间距约 $1.46\fm$；
        $n_{\max}=6$ 给出约 $-305\MeV$、对间距压到 $0.7$--$0.9\fm$。
        实验氚核束缚能 $-8.48\MeV$、半径约 $1.7\fm$。
\end{enumerate}

\begin{table}[htbp]
\centering
\caption{实验与 GEMSTORE 静态 OBE 的对照。
NN 取 \texttt{nn\_I0}（$n_{\max}=20$，$r_{\max}=30\fm$）；
NNN 取 \texttt{nnn}（同位旋重耦、$S_3$、费米子 Pauli，$n_{\max}=4$）。
$n_{\max}=6$ 为一次性扫描，数字仍列于表中，输入已从 \texttt{test/} 移除。
存档 NNN 态尚未乘 $\nu_\lambda=(4/3)\nu_\rho$；乘上后过束缚更深，见表~\ref{tab:nlam}。}
\label{tab:exp-vs-calc}
\begin{tabular}{@{}lcccc@{}}
\toprule
体系 & 量 & 实验 & 计算 & 输入 \\
\midrule
氘核 $I=0$ & $E_{\mathrm{bind}}$ & $-2.224\MeV$ & $-2.231\MeV$
  & \texttt{nn\_I0.json} \\
氘核 & 半径 & $2.13\fm$ (电荷) & $3.65\fm$ (GEM 质量 RMS)
  & $g_\omega=12.92$ \\
氚核 $I=1/2$ & $E_{\mathrm{bind}}$ & $-8.48\MeV$ & $-67.1\MeV$
  & $n_{\max}=4$, $L_{\max}=2$ \\
氚核 & 对间距 & $\sim 1.7\fm$ & $1.46$--$1.47\fm$
  & \texttt{nnn.json} \\
氚核 & $E_{\mathrm{bind}}$ & $-8.48\MeV$ & $-304.7\MeV$
  & $n_{\max}=6$, $L_{\max}=2$ \\
氚核 & 对间距 & $\sim 1.7\fm$ & $0.90/0.72/0.72\fm$
  & 扫描（未保留） \\
\bottomrule
\end{tabular}
\end{table}

需要回答的四个问题是：

\begin{enumerate}
  \item 过束缚来自三体 GEM 算法接错（动能、自旋耦合、同位旋、反对称），还是来自势本身？
  \item 若来自势，``短程排斥不够''具体指哪一段径向区间、哪一个交换？
  \item 是否必须引入自旋--轨道耦合 $V_{LS}$ 才能救氚核？
  \item Liu 的分子 OBE 是否适合拿来算三体分子态，以及是否适合拿核三体当验收？
\end{enumerate}

%----------------------------------------------------------------------
\section{物理模型与计算设置}
%----------------------------------------------------------------------

\subsection{静态弹性 OBE}

势模块是 \texttt{MODEL\_OBESTAT}，对应 OBELib 中
\texttt{OBEPotential} 取 \texttt{StaticBaryon -> True} 的极限：
Dirac 质量 $\to\infty$，没有 Gordon 混合 $g+f$，没有 $q^2/8m_B^2$，
\begin{equation}
  V_{LS}\equiv 0.
\end{equation}
弹性静态势只保留
\begin{equation}
  V(r)=V_C(r)+V_{ss}(r)\,(\vec s_i\cdot\vec s_j)+V_t(r)\,S_{12},
\end{equation}
其中 Dirac 腿上 $\vec s=\vec\sigma$，矢量介子腿上 $\vec s=\vec S$。
交换粒子固定为 $\pi,\sigma,\rho,\omega$（质量、$G$ 宇称、HeavySign 写在
\texttt{obemodel.c} 内部表）。
弹性 $0^-$ 的 $\pi$ 交换恒为零；$P\leftrightarrow P^*$ 跃迁尚未实现。
代码中 \texttt{OBEVls1}、\texttt{OBEVls2} 直接返回 $0$。

核子参数取 Liu 组 SU(3) 核子顶点（由 $\Xi$ 反演），再单独把 $g_\omega$
按氘核束缚能微调：
\begin{equation}
\begin{aligned}
  g_\pi&=0.89,\quad
  g_\sigma=10.2,\quad
  g_\rho=2.9,\quad
  g_\omega=12.92,\\
  \kappa_\rho&=3.7,\quad
  \kappa_\omega=-0.12,\quad
  M=0.940\GeV,\quad
  \Lambda=1.10\GeV.
\end{aligned}
\end{equation}
$f_\rho=\kappa_\rho g_\rho=10.73$，$f_\omega=\kappa_\omega g_\omega=-1.5504$。
这不是 CD-Bonn / AV18 一类核力拟合，而是分子 OBE 的核子尺度移植。

同位旋因子对两个 $t=1/2$ 腿：
\begin{equation}
  \langle\vec\tau_i\cdot\vec\tau_j\rangle
  =2t_{ij}(t_{ij}+1)-2t_i(t_i+1)-2t_j(t_j+1)
  =
  \begin{cases}
    -3, & I_{\mathrm{pair}}=0,\\
    +1, & I_{\mathrm{pair}}=1.
  \end{cases}
\end{equation}
$\sigma$、$\omega$ 始终取同位旋因子 $+1$。

\subsection{两体 GEM（氘核）}

\texttt{mol2\_obestat\_GEM} 按 $J$ 自动生成 $L$ 与 $L+2$ 通道。
氘核：$I=0$（\texttt{iso\_pi=iso\_rho=-3}），$J=1$，$L=0$，
通道为 ${}^3S_1$--${}^3D_1$。
动能
\begin{equation}
  T=m_1+m_2+\frac{p^2}{2\mu},\qquad \mu=\frac{m_1 m_2}{m_1+m_2}.
\end{equation}
Dirac 腿上 $V_{ss}$、$V_t$ 的自旋算符相对 $\vec s=\vec\sigma/2$ 再乘 $k_{ss}=2$，
使 $\vec\sigma\cdot\vec\sigma$ 的矩阵元与 Liu 表一致。
广义本征值 $Hc=ENc$，输出质量与 $E_{\mathrm{bind}}=M-(m_1+m_2)$。

\subsection{\texorpdfstring{三体 Jacobi GEM（氚核 / $DKK^*$）}{三体 Jacobi GEM（氚核 / DKK*）}}

三体哈密顿量为非相对论形式
\begin{equation}
\label{eq:H3}
  H=\sum_{i=1}^3 m_i
  +\sum_{i=1}^3 \frac{\vec p_i^{\,2}}{2m_i}
  +\sum_{p=12,23,31}\bigl(V_C^{(p)}+V_{ss}^{(p)}+V_t^{(p)}\bigr),
\end{equation}
没有 $\beta V\beta$，没有 $\sqrt{p^2+m^2}$。
径向积分走 SCDK（与重子 Jacobi GEM 同一套
\texttt{tir\_cent} / \texttt{tpi\_cent} / \texttt{t1r\_tens} /
\texttt{inteVcenPartA}），
势径向核由 \texttt{integral\_exp\_rn(OBEVc/OBEVss/OBEVtens)} 接到两体回调。

未加权 Jacobi 坐标
\begin{equation}
  \vec\rho=\vec r_1-\vec r_2,\qquad
  \vec\lambda=\frac{\vec r_1+\vec r_2}{2}-\vec r_3,
\end{equation}
约化质量 $\mu_\rho=m/2$，$\mu_\lambda=2m/3$（等质量）。
高斯基 $\exp(-\nu_\rho\rho^2-\nu_\lambda\lambda^2)$ 在三个对上置换对称时要求
\begin{equation}
\label{eq:nulam}
  \frac{\nu_\lambda}{\nu_\rho}=\frac{\mu_\lambda}{\mu_\rho}=\frac{4}{3},
\end{equation}
否则 $n_\rho=n_\lambda$ 并不给出 $\langle r_{12}^2\rangle=\langle r_{13}^2\rangle=\langle r_{23}^2\rangle$。
\texttt{mol3\_basis\_build} 在等质量 $S_3$ 情形于 \texttt{bases\_3b\_push\_full}
之后把 $\nu_\lambda$ 乘以 $4/3$。
可区分强子（$DKK^*$）与重子路径 \texttt{baryon\_basis\_build} 不乘这一因子。

总同位旋写入 JSON 键 \texttt{"I"}（C 字段为 \texttt{iso\_I}，避免
\texttt{complex.h} 把 \texttt{I} 定义成 \texttt{\_Complex\_I}）。
对 $\pi/\rho$ 做 $m$ 方案 CG 重耦
\begin{equation}
  V_p=\sum_k \langle f|p,t_k\rangle\, V\bigl(\mathrm{iso}(t_k)\bigr)\,
  \langle p,t_k|i\rangle,
\end{equation}
重叠、动能、RMS 再乘同位旋重叠。
费米子 Pauli
\begin{equation}
  (-1)^{\ell_\rho+s_{ij}+t_{ij}}=-1.
\end{equation}
三个全同 Dirac 腿且 \texttt{sym12=-1} 时只保留一张 Jacobi 图（$c_{\max}=1$）。
\texttt{system.I} 为三体分子态必填，对 \(t_{ij}\) 重耦到总同位旋；
\texttt{pair*.iso\_pi}/\texttt{iso\_rho} 被解析但不再用于势。

重叠按 $10^{-8}$ 剪枝后走标准本征值（\texttt{eigen\_standard}），
本征值是质量。
$E_{\mathrm{bind}}=M-(m_1+m_2+m_3)$，并输出三对 RMS。

\subsection{验收约定}

束缚判据：$E_{\mathrm{bind}}<0$ 且 RMS 明显小于盒子 $r_{\max}$。
近阈、RMS 贴着 $r_{\max}$ 的态视为未束缚的盒子态。
\texttt{nmax} 是每条 Jacobi 轴上几何高斯宽度的个数，不是径向激发量子数。

%----------------------------------------------------------------------
\section{\texorpdfstring{两体侧：氘核把 $g_\omega$ 钉在尾部}{两体侧：氘核把 g\_omega 钉在尾部}}
%----------------------------------------------------------------------

\subsection{耦合扫描}

SU(3) 表值 $g_\omega=8.7$ 给出氘核 $E_{\mathrm{bind}}=-21.43\MeV$、$\mathrm{RMS}=1.50\fm$，过束缚。
Bonn 量级 $g_\omega=15.85$ 在本静态、$V_{LS}=0$ 的算符上使氘核解束缚。
扫描后取 $g_\omega=12.92$，
\begin{equation}
  E_d=-2.231\MeV,\qquad \mathrm{RMS}=3.65\fm.
\end{equation}
束缚能与实验一致；质量 RMS 偏大，原因是 $\pi$ 张量顶点是 Liu 的 $g_\pi=0.89$，
不是核力常用的 $g_{\pi NN}\sim 13$，且没有 $V_{LS}$。
氘核波函数因此更晕。

\subsection{\texorpdfstring{$I=0$ ${}^3S_1$ 径向井}{I=0 3S1 径向井}}

对角高斯（$n_{\max}=1$）在氘核通道上看到的有效井大致为
（单位 $\MeV$，半径为高斯峰位）：

\begin{center}
\begin{tabular}{@{}ccccc@{}}
\toprule
$r$ ($\fm$) & $0.2$ & $0.5$ & $1.0$ & $2.0$ / $3.0$ \\
\midrule
有效 $V$ & $+214$ & $-91$ & $-40$ & $-2.4$ / $-0.3$ \\
\bottomrule
\end{tabular}
\end{center}

$r=0.2\fm$ 的排斥几乎全部来自 $V_{ss}$；$V_C$ 在那里约 $-337\MeV$，
被自旋--自旋芯顶起来。
\textbf{势并没有塌到原点}：短程芯是存在的。
吸引最深的一段在 $0.5$--$1\fm$，来自 $\sigma$ 吸引与 $\omega$ 排斥抵消后的剩余。
$r\gtrsim 2\fm$ 只剩几个 $\MeV$ 的浅尾，正是晕氘核主要采样的区域。

因此：$g_\omega$ 是被 $3.65\fm$ 处的尾约束的。
$\sigma$--$\omega$ 在 $0.9\fm$ 的精细抵消几乎不受氘核束缚能约束。

%----------------------------------------------------------------------
\section{三体侧：过束缚的定量表现}
%----------------------------------------------------------------------

\subsection{生产设置下的氚核}

生产输入：三个 Dirac 核子，$m_i=0.938\GeV$，$J^P=1/2^+$，$I=1/2$，
\texttt{sym12=-1}，$L_{\max}=2$，同一套 \texttt{param\_OBESTAT\_nn.json}。
阈能 $3\times 0.938=2.814\GeV$。

\begin{table}[htbp]
\centering
\caption{氚核通道基态（存档态，尚未乘 $\nu_\lambda=4/3\,\nu_\rho$）。}
\label{tab:nnn-prod}
\begin{tabular}{@{}cccccc@{}}
\toprule
$n_{\max}$ & $r_{\min},r_{\max}$ ($\fm$) & $E_{\mathrm{bind}}$ ($\MeV$)
  & $r_{12}$ ($\fm$) & $r_{13}$ ($\fm$) & $r_{23}$ ($\fm$) \\
\midrule
$4$ & $0.1$, $10$ & $-67.1$ & $1.47$ & $1.46$ & $1.46$ \\
$6$ & $0.1$, $10$ & $-304.7$ & $0.90$ & $0.72$ & $0.72$ \\
\bottomrule
\end{tabular}
\end{table}

$n_{\max}=6$ 的第一激发约为 $-24\MeV$，半径 $2.1$--$2.4\fm$，
仍深于实验基态。
基态半径随基组变大而变小、能量变深，说明变分原理在挖一个中程口袋，
而不是数值鬼态（$c^T N c\sim 1$）。

\subsection{\texorpdfstring{等质量 $S_3$ 的 $\nu_\lambda$ 比例}{等质量 S3 的 nu\_lambda 比例}}

未加权 Jacobi 上 $\nu_\rho=\nu_\lambda$ 时，单高斯
$\langle r_{13}^2\rangle/\langle r_{12}^2\rangle=\sqrt{5/4}\approx 1.12$。
乘式~\eqref{eq:nulam} 后，$n_{\max}=1$ 的三对 RMS 相等。
表~\ref{tab:nlam} 给出乘完之后的结果。

\begin{table}[htbp]
\centering
\caption{等质量 $S_3$ 取 $\nu_\lambda=(4/3)\nu_\rho$ 之后的 NNN。
对称 S 波地板已经是 $-76\MeV$；$L_{\max}=2$ 再挖到 $-300\MeV$ 量级。}
\label{tab:nlam}
\begin{tabular}{@{}lccc@{}}
\toprule
设置 & $E_{\mathrm{bind}}$ ($\MeV$) & 对间距 ($\fm$) & 说明 \\
\midrule
$L_{\max}=0$，$n_{\max}=1$，$r=1.0\fm$ & $-76$ & $0.87$ (三对相等) & S 波口袋 \\
$L_{\max}=0$，$n_{\max}=6$ & $-119$ & $r_{12}=1.03$, $r_{13}=0.87$ & 多宽度再挖 \\
$L_{\max}=2$，$n_{\max}=4$ & $-83$ & --- & D 波混合 \\
$L_{\max}=2$，$n_{\max}=6$ & $-316$ & $0.87/0.69/0.69$ & 生产基组 \\
\bottomrule
\end{tabular}
\end{table}

$n_\rho\neq n_\lambda$ 只改变变分空间的几何覆盖，不改变口袋深度的量级。
乘 $4/3$ 之后过束缚略深，说明原先略不对称的基组并没有人为制造过束缚。

\subsection{Tjon 线}

经验上，氘核束缚能约 $-2.2\MeV$ 的现实核力给出氚核约 $-8\MeV$，
二者沿 Tjon 线关联~\cite{Tjon1975}。
一套在两体上只有 $-2.2\MeV$ 的势，若三体落到 $-300\MeV$，
要么两体芯/中程口袋对三体的权重与两体完全不同，
要么三体哈密顿量接错。
下一节把第二种可能逐项关掉。

%----------------------------------------------------------------------
\section{分析过程：逐项排除算法错误}
%----------------------------------------------------------------------

分析顺序按``最容易把能量整体平移或整体放大的算符''排列。
每一项都用 $n_{\max}=1$ 或解析极限做独立核对，避免被大基组变分掩盖。

\subsection{动能}

非相对论单粒子动能矩阵元
\begin{equation}
  \langle p_c^2/(2m_c)\rangle
  =\texttt{InteExprn}(b_{11},n+2)/(2m_c),
\end{equation}
其中 \texttt{tpi\_cent} 把 SCDK 变量变到动量空间，\texttt{inteRMS} 升两阶。
等质量、单高斯、$\ell_\rho=\ell_\lambda=0$ 的解析式为
\begin{equation}
  T=\frac{3\nu_\rho}{m}+\frac{9\nu_\lambda}{4m}.
\end{equation}
取 $\nu=0.009734\GeV^2$（对应 $r\sim 2\fm$ 量级）时，
数值 $T=54.5\MeV$，与解析式在 $0.1\MeV$ 内一致；
分项 $T_1=T_2=19.5\MeV$，$T_3=15.6\MeV$，符合 Jacobi 第三粒子的约化质量。
动能没有整体少乘因子或用错质量。

\subsection{重叠与质量阈值}

重叠走 \texttt{inteNfi}。对角 $N=1$。
哈密顿装配为
\begin{equation}
  H = \Bigl(\sum_i m_i\Bigr) N
  +\sum_{p=1}^{3}\bigl(T_p+V_C^{(p)}+V_{ss}^{(p)}+V_t^{(p)}\bigr),
\end{equation}
见 \texttt{cmol3.c} 中 \texttt{add\_scaled(..., mtot)} 与三对
\texttt{add\_reduced}。
阈值 $2.814\GeV$ 与 $3\times 0.938$ 一致，$E_{\mathrm{bind}}$ 没有用错减数。

\subsection{自旋--自旋与张量}

三体 $V_{ss}$、$V_t$ 用 \texttt{vss\_hadron} / \texttt{vtens\_hadron}，
对自旋取自基组，Dirac 腿同样 $k_{ss}=2$。
SCDK 在调用前按 \texttt{channel\_3b\_ijk} 置换 $m_z$。

$n_{\max}=1$、$L_{\max}=0$、$r=2\fm$ 的核对：
\begin{itemize}
  \item 纯 S 波上张量 $V_t=0$（$S_{12}$ 在 $\ell=0$ 对角为零）。
  \item 对 $12$ 的 $V_{ss}$ 在 $s_{12}$ 上对角：
        ${}^1S$（$s_{ij}=0,t_{ij}=1$）与 ${}^3S$（$s_{ij}=1,t_{ij}=0$）
        各自有非零对角元，交叉为零。
  \item 对 $13$、$23$ 在 ${}^1S$--${}^3S$ 之间有非对角元，
        数值约 $-5.5\MeV$，正是自旋 CG 重耦所要求的混合。
\end{itemize}
自旋 Clebsch--Gordan 没有把三对 $V_{ss}$ 同时放到错误的吸引通道上。
过束缚不是自旋接线事故。

\subsection{同位旋重耦}

$t_1=t_2=t_3=1/2$、$I=1/2$ 的单位性检验：
\begin{equation}
\begin{aligned}
  \langle c=1,t_{12}=0\,|\,c=1,t_{12}=0\rangle &= 1,\\
  \langle c=1,0\,|\,c=1,1\rangle &= 0,\\
  \langle c=1,t_{12}=0\,|\,c=3,t_{31}=0\rangle &= -1/2,\\
  \langle c=1,0\,|\,c=3,1\rangle &= -\sqrt{3}/2,
\end{aligned}
\end{equation}
且 $|U|^2=1$。
$\pi/\rho$ 按算符对的 $t_p$ 做完备性和，
$N$、$T$、RMS 乘同位旋重叠。
聚类输入（不对总 $I$ 重耦、三对 \texttt{iso\_pi=-1}）同样过束缚，
只是三对 RMS 更不相等。
过束缚在重耦打开之前就已经存在，重耦不是根源。

\subsection{\texorpdfstring{Pauli 与 $S_3$}{Pauli 与 S3}}

\texttt{sym12=-1} 且三 Dirac 全同时：只保留 $c=1$，
并施加 $(-1)^{\ell_\rho+s_{ij}+t_{ij}}=-1$。
关闭 Pauli、三张独立图的``可区分核子''计算给出更紧的坍缩态
（曾见到 $\sim -400\MeV$、RMS $\sim 0.6\fm$），
与``反对称把态抬高''的方向一致。
反对称已经在工作，且不足以抵消口袋。

\subsection{重子 GEM 回归}

三体分子路径与重子路径共享 SCDK 与 \texttt{matrix\_productT}
（$ABA^T$，两次 $n^3$ 乘积）。
在 \texttt{cmol3.c} 中加入等质量 $\nu_\lambda$ 比例之后，
重子 \texttt{baryon\_basis\_build} 未改。
复算存档
\texttt{ccc\_sp\_3h2\_Pp\_L0\_jl1}、
\texttt{nnc\_sm\_1h2\_Pp\_L0\_jl0}
的基态质量差分别为 $\sim 4\times 10^{-12}\GeV$ 与 $\sim 4\times 10^{-10}\GeV$。
高激发最大偏差约 $0.09\MeV$，来自 \texttt{productT} 求和顺序，不是物理移动。
三体分子侧的修改没有破坏已验收的重子谱。

\subsection{S 波对称高斯扫描}

为把``口袋''从变分语言变成径向语言，限制 $L_{\max}=0$、
$\nu_\lambda=(4/3)\nu_\rho$、$n_{\max}=1$，扫描高斯半径。
结果见表~\ref{tab:swave} 与图~\ref{fig:pocket}。

\begin{table}[htbp]
\centering
\caption{等质量、置换对称、纯 S 波、单高斯的三核子能量。
$r\lesssim 0.5\fm$ 排斥，排除原点坍缩；
$r\sim 0.9\fm$ 出现 $-76\MeV$ 地板。}
\label{tab:swave}
\begin{tabular}{@{}ccc@{}}
\toprule
高斯半径 ($\fm$) & $E_{\mathrm{bind}}$ ($\MeV$) & 定性 \\
\midrule
$0.4$ & $+516$ & 芯，强排斥 \\
$0.5$ & $+160$ & 芯边缘 \\
$0.7$ & $-50$ & 进入口袋 \\
$1.0$ & $-76$ & 口袋底（$\langle r_{ij}\rangle\simeq 0.87\fm$） \\
$1.5$ & $-40$ & 口袋外侧 \\
$2.0$ & $-17$ & 仍深于氚核 \\
$3.0$ & $-1.1$ & 接近氘核尾 \\
\bottomrule
\end{tabular}
\end{table}

\begin{figure}[htbp]
\centering
\begin{tikzpicture}[x=1.55cm,y=0.018cm,>=Stealth]
  \draw[->] (0,0) -- (6.4,0) node[right] {$r$ ($\fm$)};
  \draw[->] (0,-90) -- (0,80) node[above] {$E$ 或 $V$ ($\MeV$)};
  \foreach \y in {-80,-60,-40,-20,20,40,60} {
    \draw[gray!30] (0.05,\y) -- (6.2,\y);
  }
  \draw[dotted] (0,0) -- (6.2,0);

  \draw[thick,blue!70!black]
    (0.15,70) -- (0.25,40) -- (0.45,-20) -- (0.7,-55)
    -- (1.0,-40) -- (1.5,-18) -- (2.0,-4) -- (3.0,-0.6) -- (4.5,-0.1);
  \node[blue!70!black,font=\small] at (2.7,-14) {NN ${}^3S_1$ 井（示意）};

  \draw[thick,red!70!black]
    (0.4,70) -- (0.5,40) -- (0.7,-28) -- (1.0,-43)
    -- (1.5,-22) -- (2.0,-9.5) -- (3.0,-0.6);
  \fill[red!70!black] (1.0,-43) circle (1.6pt);
  \node[red!70!black,font=\small,align=left] at (3.55,-50)
    {对称 S 波三体\\单高斯扫描};

  \draw[dashed,olive!70!black] (3.65,-90) -- (3.65,70);
  \node[olive!70!black,font=\footnotesize,rotate=90,anchor=south]
    at (3.52,10) {氘核 RMS $3.65\fm$};
  \draw[dashed,orange!80!black] (0.87,-90) -- (0.87,20);
  \node[orange!80!black,font=\footnotesize,rotate=90,anchor=south]
    at (0.74,-20) {三体口袋 $0.87\fm$};

  \foreach \x in {1,2,3,4,5} {
    \draw (\x,2) -- (\x,-2) node[below] {\x};
  }
\end{tikzpicture}
\caption{示意：氘核晕半径落在势的浅尾上，三体对称 S 波落在 $\sigma/\omega$
中程口袋里。纵轴对 NN 井与三体扫描用了不同的直观缩放，只表示相对位置，
精确数字见表~\ref{tab:swave}。}
\label{fig:pocket}
\end{figure}

这一条扫描不依赖 $L_{\max}=2$、不依赖 $n_\rho\neq n_\lambda$、
不依赖张量、不依赖 $V_{LS}$。
\textbf{只要三对同时坐进 $0.6$--$1.2\fm$ 的吸引井，三体就已经过束缚一个数量级。}

%----------------------------------------------------------------------
\section{问题所在}
%----------------------------------------------------------------------

\subsection{诊断结论}

\begin{center}
\fbox{\parbox{0.92\textwidth}{\raggedright
\textbf{主因：} 静态 Liu OBE 在 $0.6$--$1.2\fm$ 留下过深的 $\sigma/\omega$ 吸引口袋。
氘核只标定了 $r\gtrsim 2\fm$ 的尾，$g_\omega$ 不能约束该口袋。
三体 S 波三对同时采样口袋，给出 $\gtrsim 70\MeV$ 的过束缚；
更大基组与 $L_{\max}=2$ 把能量继续挖到 $-300\MeV$ 量级。\\[0.4em]
\textbf{非主因：} 动能接线、自旋 CG、同位旋 Racah、Pauli、$n_\rho\neq n_\lambda$、
重子 SCDK、缺少 $V_{LS}$（对纯 S 波）。
}}
\end{center}

``短程排斥不够''在本问题里应当说成：
\textbf{芯（$r\lesssim 0.5\fm$）已经排斥，真正不够的是中程（$0.6$--$1.2\fm$）
$\omega$ 对 $\sigma$ 的抵消}。
原点没有塌缩；$r=0.4\fm$ 的对称三体高斯是 $+516\MeV$。
变分把系统停在口袋底，而不是冲进硬芯。

\subsection{为什么两体对、三体错}

设两体相对波函数为 $\psi_d(r)$，三体对称 S 波对密度为 $\rho_3(r)$。
氘核 $\langle r\rangle\sim 3.65\fm$，
\begin{equation}
  \int_{0}^{1.2\fm} r^2|\psi_d(r)|^2\,\mathrm{d}r \ll 1,
\end{equation}
对口袋几乎没有权重。
三体要把三对同时拉近，对密度在 $0.9\fm$ 附近有最大权重，
有效势大约是三对口袋之和减去多体动能。
单对口袋深 $\sim 40$--$90\MeV$ 时，三对相加很容易超过 $100\MeV$，
动能只抵消一部分，剩下 $-76\MeV$ 的 S 波地板。

这与核力现象学一致：现实 NN 势在中程必须被精细平衡，
并且通常还需要三体力才能同时描述 $d$ 与 ${}^3\mathrm{H}$。
本模型用一个近阈分子 OBE 的 $g_\omega$ 去兼做这两件事，几何上做不到。

\subsection{\texorpdfstring{$V_{LS}$ 救不了 S 波地板}{V\_LS 救不了 S 波地板}}

静态极限把 $V_{LS}$ 设为零是模型定义，不是实现遗漏。
即便恢复有限质量 OBE 的 $V_{LS}$，
在 $L=0$ 的矩阵元
\begin{equation}
  \langle \ell=0\,|\,\vec L\cdot\vec S\,|\,\ell=0\rangle=0.
\end{equation}
表~\ref{tab:swave} 的 $-76\MeV$ 地板全部来自 $V_C+V_{ss}$。
$L_{\max}=2$ 之后 $V_{LS}$ 可以移动 D 波混合，
但不能把已经过深的 S 波口袋抬回 $-8\MeV$。
因此``需要引入自旋--轨道耦合才能解决 NNN''这一判断不成立。

\subsection{Liu OBE 的适用域}

Liu 2018~\cite{Liu2018} 的静态 OBE 是为近阈两体分子
（如 $D^*\Xi$）设计的：束缚能 $\mathrm{keV}$--$\mathrm{MeV}$，半径数 $\fm$ 到盒子尺度。
它的耦合来自 SU(3) / VMD / 重夸克顶点，不是相移与氘核、氚核的联合拟合。

\begin{itemize}
  \item \textbf{适合：} 松散三体分子，尤其每对都带一条 $0^-$ 腿、
        因而 $V_{ss}=V_t=0$、只剩较弱 $V_C$ 的系统。
        $DKK^*$ 正是这一类：两体 DK、DK$^*$、KK$^*$ 均未束缚
        （$E_{\mathrm{bind}}\sim +0.1\MeV$，RMS $\sim 35\fm$），
        三体 $1^-$ 随 $r_{\max}$ 走盒子，
        $E_{\mathrm{bind}}r_{\max}^2\sim 210\MeV\,\fm^2$。
        NNN 过束缚并不蕴含 $DKK^*$ 的 GEM 接线错误。
  \item \textbf{不适合：} 把氚核 $-8.48\MeV$ 当作该 OBE 的验收标准，
        或把它当作核多体哈密顿量去预测核结合能。
\end{itemize}

\subsection{与电荷半径的关系}

NN 的 GEM RMS $3.65\fm$ 大于氘核电荷半径 $2.13\fm$，
在``程序输出是质量均方根''的前提下可以接受，
也与晕氘核、偏弱的 $\pi$ 张量一致。
这一条只说明两体尾偏软，
\textbf{不能}反过来证明中程口袋已经被氘核约束好。
NNN 半径在 $n_{\max}=4$ 时看起来接近 $1.7\fm$，
只是变分还没挖到底；$n_{\max}=6$ 立刻掉到 $0.7\fm$ 并过束缚到 $-300\MeV$。

%----------------------------------------------------------------------
\section{拟解决方案}
%----------------------------------------------------------------------

按目标分成三条互不替代的路径。
\textbf{不要}用氚核束缚能去回标 $g_\omega$：
加大 $g_\omega$ 会先解束缚氘核（扫描中 $g_\omega=15.85$ 已解束缚），
这正是 Tjon 线所禁止的单参数挪动。

\subsection{路径 A：验收三体 GEM 本身（优先）}

NNN 不能当 GEM 验收，因为势的中程口袋已经错了一个数量级。
应换用答案已知的示意图势，只测算法：

\begin{enumerate}
  \item 中心高斯吸引阱或 MT-V / Yamaguchi 可分势，
        两体束缚能与三体基态有文献或独立数值。
  \item 等质量三玻色子、再加费米子 Pauli，分别核对
        动能、重叠、置换、同位旋重耦。
  \item 纯 S 波解析极限（单高斯、式~\eqref{eq:nulam}）作为持续回归。
\end{enumerate}

通过之后，三体 GEM 视为可信；$DKK^*$ 的未束缚结论可以保留。

\subsection{\texorpdfstring{路径 B：坚持分子物理（$DKK^*$ 等）}{路径 B：坚持分子物理（DKK* 等）}}

保持当前静态弹性 OBE，不引入 $P\leftrightarrow P^*$，除非另有物理需求。
\begin{itemize}
  \item 验收标准改为：两体分子谱与近阈行为、三体是否相对两体阈值束缚、
        RMS 是否离开盒子。
  \item 不再用 $d$ / ${}^3\mathrm{H}$ 的 MeV 级核数据卡分子 OBE。
  \item NNN 作为``该势移植到核子上会过束缚''的警示留下，不当作程序缺陷。
\end{itemize}

若将来打开 $P\leftrightarrow P^*$，那是额外的分子动力学通道，
仍不是核力拟合。

\subsection{路径 C：若目标改为核三体}

必须换势，而不是在 Liu 顶点上继续拧 $g_\omega$。

\begin{enumerate}
  \item 采用核力 NN（AV18、CD-Bonn、或至少含有被相移约束的中程 $\sigma/\omega$
        与接触项的 OBE），再加不可约三体力（Fujita--Miyazawa / $\mathrm{N}^2\mathrm{LO}$ 3NF）。
  \item $V_{LS}$、张量、P 波相移一起进入拟合；氚核是拟合目标之一，不是事后检验。
  \item 这一路径与分子模块分开：核力势不应覆盖 \texttt{param\_OBESTAT\_dk*.json}。
\end{enumerate}

\subsection{明确不采取的步骤}

\begin{itemize}
  \item 不把 $g_\omega$ 调到使氚核到 $-8.48\MeV$。
        该操作会毁掉已经对齐的氘核，并掩盖口袋问题。
  \item 不把 $V_{LS}$ 当作 S 波过束缚的补丁。
  \item 不把 $n_\rho=n_\lambda$ 过滤器当作物理修正；
        等质量 $S_3$ 只保留 $\nu_\lambda=(4/3)\nu_\rho$。
  \item 不改重子 \texttt{baryon\_basis\_build} 的宽度约定。
  \item 不为了 NNN 去改 $DKK^*$ 的聚类同位旋路径。
\end{itemize}

\subsection{建议的执行顺序}

\begin{enumerate}
  \item 用示意图势做三体 GEM 的独立验收（路径 A）。
  \item 在文档与测试说明中写明：Liu OBE + 氚核不是验收对。
  \item 分子谱继续走路径 B。
  \item 仅当明确要做核三体时，另开核力势模块（路径 C）。
\end{enumerate}

%----------------------------------------------------------------------
\section{\texorpdfstring{对 $DKK^*$ 与现有代码的含义}{对 DKK* 与现有代码的含义}}
%----------------------------------------------------------------------

$DKK^*$ 的三对都含 $0^-$ 腿，弹性静态势上 $V_{ss}=V_t=0$，
哈密顿量是阈值质量 + 非相对论动能 + 三对中心 $V_C$。
两体 DK / DK$^*$ / KK$^*$ 均未束缚，三体随盒子变大而趋于阈值，
与``弱吸引、无分子束缚''一致。
NNN 口袋来自 Dirac--Dirac 的 $V_{ss}$ 与强 $\sigma/\omega$，
在 $DKK^*$ 上不出现同一机制。

代码层面应保持：
\begin{itemize}
  \item \texttt{cmol3.c} 等质量 $S_3$ 的 $\nu_\lambda=(4/3)\nu_\rho$；
  \item \texttt{system.I} 必填，始终按总同位旋重耦；
  \item 重子路径不动；
  \item SPECTRA 输出仍放在 \texttt{test/Spectra-Molecule2/} 与
        \texttt{test/Spectra-Molecule3/}。
\end{itemize}

%----------------------------------------------------------------------
\section{结论}
%----------------------------------------------------------------------

NNN 束缚能对不上，是当前分子 OBE 移植到核子上最严重的物理问题，
但不是三体 GEM 算符接错。
动能、自旋、同位旋、Pauli、重子回归均通过；
对称 S 波单高斯已经给出 $-76\MeV$ 的地板。
问题坐落在 $0.6$--$1.2\fm$ 的 $\sigma/\omega$ 口袋：
氘核晕态标定了尾，$g_\omega$ 约束不到口袋，
三体三对同时落入口袋后沿 Tjon 线严重偏离。
$V_{LS}$ 在 $L=0$ 上为零，不是补救手段。
Liu 静态 OBE 适合松散三体分子，不适合当核三体哈密顿量。
下一步应以示意图势验收 GEM，并把氚核实验值从分子 OBE 的验收清单中移除。

%----------------------------------------------------------------------
\section*{附录 A  \ 关键输入与代码位置}
%----------------------------------------------------------------------

\begin{table}[htbp]
\centering
\caption{本报告引用的输入与驱动。}
\label{tab:files}
{\small
\begin{tabular}{@{}p{0.62\textwidth}p{0.30\textwidth}@{}}
\toprule
路径 & 内容 \\
\midrule
\texttt{src/model/obemodel.c} & 静态 $\pi\sigma\rho\omega$，$V_{LS}\equiv 0$ \\
\texttt{src/model/cmol2.c} & 两体 GEM，$T=m_1+m_2+p^2/2\mu$ \\
\texttt{src/model/cmol3.c} & 三体 Jacobi GEM，式~\eqref{eq:H3} \\
\texttt{test/Spectra-Molecule2/param\_OBESTAT\_nn.json} & $g_\omega=12.92$ \\
\texttt{test/Spectra-Molecule2/nn\_I0.json} & 氘核 $I=0$，$n_{\max}=20$ \\
\texttt{test/Spectra-Molecule3/nnn.json} & 氚核 $n_{\max}=4$，$I=1/2$ \\
\texttt{test/Spectra-Molecule3/dkkstar.json} & $DKK^*$ 聚类，无 \texttt{I} \\
\bottomrule
\end{tabular}
}
\end{table}

CLI 必须使用 \texttt{gemstore --compute file.json}。
仅给出 JSON 文件名时程序只打印标识后退出。

%----------------------------------------------------------------------
\section*{附录 B  \ 本征值与半径的原始数字}
%----------------------------------------------------------------------

氘核（\texttt{nn\_I0.state.json}）：
\begin{align}
  M&=1.873769259\GeV, &
  m_1+m_2&=1.876\GeV,\\
  E_{\mathrm{bind}}&=-2.231\MeV, &
  \mathrm{RMS}&=3.649\fm.
\end{align}

氚核 $n_{\max}=4$（\texttt{nnn.state.json}）：
\begin{equation}
  M=2.746933\GeV,\quad
  E_{\mathrm{bind}}=-67.067\MeV,\quad
  r_{12}=1.472\fm,\;
  r_{13}=r_{23}=1.457\fm.
\end{equation}

氚核 $n_{\max}=6$（扫描，输入未保留）：
\begin{equation}
  M=2.509336\GeV,\quad
  E_{\mathrm{bind}}=-304.664\MeV,\quad
  r_{12}=0.904\fm,\;
  r_{13}=r_{23}=0.715\fm.
\end{equation}

\begin{thebibliography}{9}

\bibitem{Liu2018}
M.-Z.~Liu, Y.-W.~Pan, F.-Z.~Peng, M.~S\'anchez S\'anchez, L.-S.~Geng,
A.~Hosaka, and M.~Pav\'on Valderrama,
Phys.\ Rev.\ D \textbf{98}, 014014 (2018),
\href{https://arxiv.org/abs/1805.08384}{arXiv:1805.08384}.

\bibitem{Tjon1975}
J.~A.~Tjon,
Phys.\ Lett.\ B \textbf{56}, 217 (1975).

\end{thebibliography}

\end{document}
